% block: matrix-fractional-power-heading page:220 type:section
\subsection{{矩阵分数次幂}}\label{sec:matrix-fractional-power}\label{srcpage:printed-208}

% block: matrix-fractional-power-definition page:220 type:ordinary-paragraph
{矩阵分数次幂可利用 $\log(A)$ 如下定义:}
\begin{equation}\label{eq:matrix-fractional-power-definition}A^\alpha=\exp(\alpha\log(A)),\end{equation}
{其中 $0<\alpha<1$. 条件 $0<\alpha<1$ 看似限制过强, 但在实际中已经足够. 事实上, 如果想计算 $A^{2.3}$, 则 $A^{2.3}$ 可分解为 $A^{2.3}=A^2A^{0.3}$. 因此, 问题在于如何计算 $A^{0.3}$, 它对应于 \eqref{eq:matrix-fractional-power-definition} 中 $\alpha=0.3$ 的情形.}

% block: matrix-fractional-power-integral page:220 type:ordinary-paragraph
{通过 \eqref{eq:matrix-fractional-power-definition} 中的定义计算 $A^\alpha$ 需要两个矩阵函数: 矩阵对数和矩阵指数函数. 另一方面, 下列积分形式中不出现这两个矩阵函数:}
\begin{equation}\label{eq:matrix-fractional-power-integral}A^\alpha=\frac{\sin(\alpha\pi)}{\alpha\pi}A\int_0^\infty(t^{1/\alpha}I+A)^{-1}\,\mathrm{d}t\quad(0<\alpha<1).\end{equation}

% block: matrix-fractional-power-de-formula page:220 type:ordinary-paragraph
{文献 \BookCite{186} 考虑将 DE 公式应用于 \eqref{eq:matrix-fractional-power-integral}, 如下所述:}
\begin{equation}\label{eq:matrix-fractional-power-de-integral}A^\alpha=\int_{-\infty}^\infty F_{\mathrm{DE}}(x)\,\mathrm{d}x,\end{equation}
{其中}
\begin{equation}\label{eq:matrix-fractional-power-de-integrand}F_{\mathrm{DE}}(x)=t'(x)F(t(x)),\quad F(t)=\frac{\sin(\alpha\pi)}{\alpha\pi}A(t^{1/\alpha}I+A)^{-1}.\end{equation}

% block: matrix-fractional-power-algorithm-transition page:220 type:ordinary-paragraph
{利用 DE 公式计算 $A^\alpha$ 的算法见 Algorithm~\ref{alg:matrix-fractional-power-de}. Algorithm~\ref{alg:matrix-fractional-power-de} 的理论细节见 \BookCite{186}.}

% block: matrix-fractional-power-de-algorithm page:220 type:ordinary-paragraph
\begin{algorithm}[H]\caption{{计算 $A^\alpha$ 的 $m$ 点 DE 公式}}\label{alg:matrix-fractional-power-de}\begin{algorithmic}[1]
\State \textbf{Input} $A\in\mathbb{R}^{n\times n}$, $\alpha\in(0,1)$, $\epsilon>0$, $m$
\State $l,r=\mathtt{GetInterval}(A,\alpha,\epsilon)$
\State {设置} $\widetilde{F}_{\mathrm{DE}}(x):=\exp(\alpha\pi\sinh(x)/2)\cosh(x)[\exp(\pi\sinh(x)/2)I+A]^{-1}$.
\State $h=(r-l)/(m-1)$
\State $T=h[\widetilde{F}_{\mathrm{DE}}(l)+\widetilde{F}_{\mathrm{DE}}(r)]/2+h\sum_{k=1}^{m-2}\widetilde{F}_{\mathrm{DE}}(l+kh)$
\State \textbf{Output} $\sin(\alpha\pi)AT/2\approx A^\alpha$
\State{}
\State \textbf{function} $\mathtt{GetInterval}(A,\alpha,\epsilon)$
\State {计算 $\|A\|$, $\|A^{-1}\|$.}
\State $a_1=[\alpha\pi(1+\alpha)\epsilon]/[4\sin(\alpha\pi)(1+2\alpha)]$, $a_2=(2\|A^{-1}\|)^{-\alpha}$
\State $a=\min\{a_1,a_2\}$
\State $\begin{aligned}b_1&=[\pi(1-\alpha)(2-\alpha)\epsilon]^{\alpha/(\alpha-1)}/[4\sin(\alpha\pi)(3-2\alpha)\|A\|]^{\alpha/(\alpha-1)},\\b_2&=(2\|A\|)^\alpha\end{aligned}$
\State $b=\max\{b_1,b_2\}$
\State $l=\operatorname{asinh}(2\log(a)/\alpha\pi)$, $r=\operatorname{asinh}(2\log(b)/\alpha\pi)$
\State \textbf{return} $l,r$
\State \textbf{end function}
\end{algorithmic}\end{algorithm}

% block: matrix-fractional-power-column-systems page:221 type:ordinary-paragraph
{由 Algorithm~\ref{alg:matrix-fractional-power-de}, 算得的 $A^\alpha$ 的第 $j$ 列为 $(\sin(\alpha\pi)AT/2)\boldsymbol{e}_j$. 因此, 由第 5 行, 必须利用下列方程计算各 $\widetilde{F}_{\mathrm{DE}}(l+ih)\boldsymbol{e}_j$:}
\[e^{\alpha\pi\sinh(l+ih)/2}\cosh(l+ih)(e^{\pi\sinh(l+ih)/2}I+A)\boldsymbol{x}^{(i)}=\boldsymbol{e}_j\quad\text{for }i=0,1,\ldots,m-1.\]
{类似于 \eqref{eq:matrix-log-quadrature-systems} 和 \eqref{eq:matrix-log-shifted-systems}, 上述方程可改写为移位线性方程组. 因此, 如果矩阵 $A$ 是大型稀疏矩阵, 使用 (移位) Krylov 子空间方法将是一种选择.}
